Sparse co-occurrence between AI functions and decision functions is easily read as a structural hole between artificial intelligence (AI) and managerial decision-making. This project drew that inference from its own literature matrix and built a multi-level programming model on it. We show why co-occurrence cannot carry the structural-hole construct and test the premise with flow data. Study~1 re-examines the literature programme: 16 depleted function pairs keep their direction under every check we could run, but lose significance under realistic misclassification; the direction is descriptive and possibly an artefact of non-random coding errors. Study~2 applies four preregistered analyses to 8,512,829 US utility patents. AI and decision patents cite each other far above a time-stratified expectation (OE 3.05, 95\% CI 2.94--3.17), and both relative hole criteria fail in all 6 sensitivity checks. Specialist AI and decision firms remain separated beyond a portfolio-preserving null ($p = .001$), but the residual is small: the null reproduces 85\% to 91\% of the separation. A brokerage advantage for mixed-portfolio firms cannot be established. Ruling out a knowledge-flow hole does not establish an organisational bottleneck. That account is one candidate among competing explanations, which a planned enterprise study is designed to discriminate.
